% Luc Destombe --- licence CC BY-NC-SA 4.0
% https://creativecommons.org/licenses/by-nc-sa/4.0/deed.fr
% Fichier autonome : compiler avec pdflatex, deux fois (signets, nombre de pages).
\def\lyceeversion{1.0}
\documentclass[a4paper,10pt]{article}
\makeatletter
\newif\iflycee@niveaux
\lycee@niveauxtrue
\newif\iflycee@cm
\newif\iflycee@palatino

\RequirePackage[utf8]{inputenc}
\RequirePackage[T1]{fontenc}
\RequirePackage[french]{babel}
\RequirePackage{etoolbox}

\newcommand{\ShorthandsOff}{\@ifundefined{shorthandoff}{}{\shorthandoff{;:!?}}}
\newcommand{\ShorthandsOn}{\@ifundefined{shorthandon}{}{\shorthandon{;:!?}}}
\ShorthandsOff
\AtBeginDocument{\ShorthandsOn}
\AtBeginEnvironment{tikzpicture}{\ShorthandsOff}

\RequirePackage{geometry}
\RequirePackage{fancyhdr}
\RequirePackage{lastpage}
\RequirePackage{multicol}
\RequirePackage{array}
\RequirePackage{enumitem}
\RequirePackage{xcolor}
\RequirePackage{needspace}

\RequirePackage{amsmath,amssymb}
\iflycee@cm\else\RequirePackage{fourier}\fi
\DeclareMathAlphabet{\mathcal}{OMS}{cmsy}{m}{n}
\renewcommand{\ttdefault}{lmtt}
\begingroup\catcode`\_=\active \catcode`\^=\active
\global\let_\sb \global\let^\sp
\endgroup
\AtBeginDocument{\catcode`\_=12 \mathcode`\_="8000
  \catcode`\^=12 \mathcode`\^="8000 }
\newlength{\retraitcalculs}
\setlength{\retraitcalculs}{2em}

\RequirePackage{tikz}
\usetikzlibrary{arrows.meta,calc,babel,shapes.misc}
\RequirePackage[most]{tcolorbox}
\tcbuselibrary{skins,breakable}

\definecolor{coulDef}{HTML}{1F4E79}
\definecolor{coulProp}{HTML}{7030A0}
\definecolor{coulRem}{HTML}{C55A11}
\definecolor{coulEx}{HTML}{2E7D32}
\definecolor{coulExo}{HTML}{B03A2E}
\definecolor{coulPart}{HTML}{1F4E79}
\definecolor{coulMeth}{HTML}{1B6E4F}
\definecolor{coulCourbe}{HTML}{0B5ED7}
\definecolor{coulCourbeB}{HTML}{C00000}
\definecolor{coulCourbeC}{HTML}{2E7D32}
\definecolor{coulGrille}{HTML}{8C8C8C}
\definecolor{coulRep}{HTML}{1B6E4F}
\definecolor{coulBar}{HTML}{2E7D32}

\newcommand{\R}{\mathbb{R}}
\newcommand{\Rs}{\mathbb{R}^{*}}
\newcommand{\N}{\mathbb{N}}
\newcommand{\Ns}{\mathbb{N}^{*}}
\newcommand{\Z}{\mathbb{Z}}
\newcommand{\Q}{\mathbb{Q}}
\renewcommand{\le}{\leqslant}
\renewcommand{\ge}{\geqslant}
\renewcommand{\approx}{\simeq}
\newcommand{\vect}[1]{\overrightarrow{#1}}
\newcommand{\norme}[1]{\left\lVert #1 \right\rVert}
\newcommand{\intervalle}[2]{\left[#1\,;#2\right]}
\RequirePackage{cancel}
\renewcommand{\CancelColor}{\color{coulExo}}
\newcommand{\fois}[1]{\times{\color{coulExo}#1}}
\newcommand{\barre}[1]{\cancel{#1}}
\newlength{\lycee@eqh}\newlength{\lycee@eqd}
\newsavebox{\lycee@eqa}\newsavebox{\lycee@eqb}\newsavebox{\lycee@eqc}
\newcommand{\lycee@eqcol}[1]{\hbox{$\renewcommand{\arraystretch}{1.35}%
  \begin{array}[t]{@{}>{\displaystyle}r@{}>{\displaystyle{}}l@{}}#1\end{array}$}}
\newcommand{\lycee@eqmesure}[1]{%
  \ifdim\ht#1>\lycee@eqh\setlength{\lycee@eqh}{\ht#1}\fi
  \ifdim\dp#1>\lycee@eqd\setlength{\lycee@eqd}{\dp#1}\fi}
\newcommand{\lycee@eqtrait}{\hspace{1.2em}\tikz[baseline=0pt]\draw[dash pattern=on 2pt off 2pt,
  line width=0.5pt] (0,\lycee@eqh) -- (0,-\lycee@eqd);\hspace{1.2em}}
\newcommand{\eqproduit}[2]{\par\smallskip\noindent C'est une équation produit :
  $A\times B=0 \iff A=0$ ou $B=0$.\par\nopagebreak\smallskip
  \begingroup\setlength{\lycee@eqh}{0pt}\setlength{\lycee@eqd}{0pt}%
  \setbox\lycee@eqa=\lycee@eqcol{#1}\setbox\lycee@eqb=\lycee@eqcol{#2}%
  \lycee@eqmesure\lycee@eqa\lycee@eqmesure\lycee@eqb%
  \noindent\hspace*{\retraitcalculs}\usebox\lycee@eqa\lycee@eqtrait\usebox\lycee@eqb\endgroup\par\smallskip}
\newcommand{\eqproduittrois}[3]{\par\smallskip\noindent C'est une équation produit :
  $A\times B\times C=0 \iff A=0$ ou $B=0$ ou $C=0$.\par\nopagebreak\smallskip
  \begingroup\setlength{\lycee@eqh}{0pt}\setlength{\lycee@eqd}{0pt}%
  \setbox\lycee@eqa=\lycee@eqcol{#1}\setbox\lycee@eqb=\lycee@eqcol{#2}\setbox\lycee@eqc=\lycee@eqcol{#3}%
  \lycee@eqmesure\lycee@eqa\lycee@eqmesure\lycee@eqb\lycee@eqmesure\lycee@eqc%
  \noindent\hspace*{\retraitcalculs}\usebox\lycee@eqa\lycee@eqtrait\usebox\lycee@eqb\lycee@eqtrait\usebox\lycee@eqc\endgroup
  \par\smallskip}

\tikzset{grille/.style={coulGrille,line width=0.4pt}}
\newcommand{\quadrillage}[4]{\draw[grille,step=1] (#1,#2) grid (#3,#4);}
\tikzset{
  croix/.style={cross out,draw,line width=0.6pt,inner sep=0pt,minimum size=4pt},
  croixdroite/.style={croix,rotate=45,line width=0.8pt},
}
\newcommand{\pt}[3]{\path (#1) node[croix]{} node[#2,font=\small,inner sep=2.5pt]{$#3$};}

\newcommand{\lycee@repere}[4]{%
  \draw[grille,step=1] (#1,#2) grid (#3,#4);
  \draw[-{Stealth[length=2mm]},thick] (#1,0)--({#3+0.4},0) node[below left=-1pt]{$x$};
  \draw[-{Stealth[length=2mm]},thick] (0,#2)--(0,{#4+0.4}) node[below left=-1pt]{$y$};
  \node[below left,font=\scriptsize,inner sep=1.5pt] at (0,0) {$0$};}
\newcommand{\lycee@gradx}[1]{\foreach \k/\l in {#1}{%
  \draw (\k,0.07)--(\k,-0.07) node[below,font=\scriptsize,inner sep=1.5pt]{$\l$};}}
\newcommand{\lycee@grady}[1]{\foreach \k/\l in {#1}{%
  \draw (0.07,\k)--(-0.07,\k) node[left,font=\scriptsize,inner sep=1.5pt]{$\l$};}}
\newcommand{\lycee@intff}[2]{\left[#1\,;#2\right]}
\newcommand{\lycee@intfo}[2]{\left[#1\,;#2\right[}
\newcommand{\lycee@intof}[2]{\left]#1\,;#2\right]}
\newcommand{\lycee@intoo}[2]{\left]#1\,;#2\right[}
\newcommand{\lycee@ens}[1]{\left\{#1\right\}}
\AtBeginDocument{%
  \providecommand{\repere}{\lycee@repere}%
  \providecommand{\gradx}{\lycee@gradx}%
  \providecommand{\grady}{\lycee@grady}%
  \providecommand{\Cf}{\mathcal{C}\sb{f}}%
  \providecommand{\Cg}{\mathcal{C}\sb{g}}%
  \providecommand{\intff}{\lycee@intff}%
  \providecommand{\intfo}{\lycee@intfo}%
  \providecommand{\intof}{\lycee@intof}%
  \providecommand{\intoo}{\lycee@intoo}%
  \providecommand{\ens}{\lycee@ens}}

\newlist{questions}{enumerate}{3}
\setlist[questions,1]{label=\arabic*\textdegree),leftmargin=*,itemsep=2pt,topsep=3pt}
\setlist[questions,2]{label=\alph*),leftmargin=*,itemsep=1pt,topsep=2pt}
\setlist[questions,3]{label=\roman*),leftmargin=*,itemsep=1pt,topsep=1pt}
\setlist[itemize]{label=\textbullet,leftmargin=*,itemsep=2pt,topsep=3pt}

\newcommand{\besoinplace}[1]{\par\penalty\@M\begingroup
  \dimen@=#1\relax\dimen@ii=\pagegoal\advance\dimen@ii-\pagetotal
  \ifdim\dimen@>\dimen@ii\ifdim\dimen@ii>\z@\vfil\fi\break\fi\endgroup}

\newcommand{\lycee@pied}{\thepage/\pageref{LastPage}}
\newcommand{\entete}[2]{%
  \pagestyle{fancy}\fancyhf{}%
  \fancyhead[L]{\footnotesize\color{coulPart}\textbf{#1}}%
  \fancyhead[R]{\footnotesize\color{coulPart}\textbf{#2}}%
  \fancyfoot[C]{\footnotesize\lycee@pied}%
  \renewcommand{\headrulewidth}{0.4pt}%
  \renewcommand{\headrule}{\hbox to\headwidth{%
    \color{coulPart!40}\leaders\hrule height \headrulewidth\hfill}}}

\RequirePackage{ccicons}
\newcommand{\lycee@licence}{{\scriptsize\color{gray}%
  \copyright\ Luc Destombe --- \ccbyncsa\ CC BY-NC-SA 4.0}}
\newif\iflycee@licence \lycee@licencetrue
\newcommand{\sanslicence}{\lycee@licencefalse}
\AtBeginDocument{\iflycee@licence\AddToHookNext{shipout/foreground}{%
  \put(0,0){\rlap{\kern\dimexpr1in+\hoffset+\oddsidemargin+\textwidth\relax
    \llap{\raisebox{-\dimexpr1in+\voffset+\topmargin+\headheight+\headsep
      +\textheight+\footskip\relax}{\lycee@licence}}}}}\fi}

\newcommand{\lycee@site}{faire-des-maths.fr}
\newcommand{\lycee@version}{{\scriptsize\color{gray}\lycee@site
  \ifdefined\lyceeversion\ --- Version~\lyceeversion\fi}}
\AtBeginDocument{\iflycee@licence\AddToHookNext{shipout/foreground}{%
  \put(0,0){\rlap{\kern\dimexpr1in+\hoffset+\oddsidemargin\relax
    \raisebox{-\dimexpr1in+\voffset+\topmargin+\headheight+\headsep
      +\textheight+\footskip\relax}{\lycee@version}}}}\fi}

\newcommand{\formulesserrees}{%
  \setlength{\abovedisplayskip}{3pt}\setlength{\belowdisplayskip}{3pt}%
  \setlength{\abovedisplayshortskip}{2pt}\setlength{\belowdisplayshortskip}{3pt}}

\setlength{\parindent}{0pt}

\RequirePackage[implicit=false,hidelinks,pdfencoding=auto,psdextra,bookmarksopen,
  bookmarksopenlevel=1,pdfcreator={},pdfproducer={}]{hyperref}
\RequirePackage{bookmark}
\hypersetup{pdfauthor={Luc Destombe},pdfkeywords={CC BY-NC-SA 4.0}}
\newcounter{lycee@signet}
\newcommand{\signet}[3][]{\stepcounter{lycee@signet}%
  \edef\lycee@dest{\if\relax\detokenize{#1}\relax
    signet.\arabic{lycee@signet}\else#1\fi}%
  \hypertarget{\lycee@dest}{}\bookmark[level=#2,dest=\lycee@dest]{#3}}
\pdfstringdefDisableCommands{%
  \def\R{R}\def\N{N}\def\Z{Z}\def\Q{Q}\def\vect#1{#1}\def\Cf{Cf}\def\Cg{Cg}\def\fois#1{×#1}%
  \def\barre#1{#1}\def\intervalle#1#2{[#1;#2]}\def\intff#1#2{[#1;#2]}%
  \def\textsuperscript#1{#1}\def\dfrac#1#2{#1/#2}\def\frac#1#2{#1/#2}%
  \def\vec#1{#1⃗}}%
\RequirePackage[official]{eurosym}

\geometry{top=0.8cm,bottom=1.3cm,left=1cm,right=1cm,footskip=0.6cm}

\pagestyle{fancy}
\fancyhf{}
\renewcommand{\headrulewidth}{0pt}
\fancyfoot[C]{\footnotesize Page \thepage{} / \pageref{LastPage}}

\newcommand{\titrefiche}[2]{%
  \begin{center}
    \Large\textbf{#1}\\[2pt]
    \large\textbf{#2}\\[2pt]
  \end{center}
  \vspace{0.10cm}}

\setlength{\columnseprule}{0.4pt}
\setlength{\columnsep}{15pt}
\renewcommand{\columnseprulecolor}{\color{black!45}}
\setlength{\multicolsep}{2pt}

\newcounter{partiefiche}
\renewcommand{\thepartiefiche}{\Roman{partiefiche}}
\newif\ifapresbandeau
\newcommand{\lycee@bandeau}[1]{%
  \par\penalty-600
  \vspace{0.08cm}%
  \noindent\colorbox{coulPart}{%
    \parbox{\dimexpr\linewidth-2\fboxsep\relax}{%
      \color{white}\bfseries\raggedright\strut#1\strut}}%
  \nopagebreak\par\nopagebreak
  \vspace{0.06cm}\nopagebreak
  \apresbandeautrue}
\newcommand{\partiefiche}[1]{%
  \refstepcounter{partiefiche}%
  \lycee@bandeau{\signet{1}{\thepartiefiche{} --- #1}\thepartiefiche{} --- #1}}
\newcommand{\partiefichesansnum}[1]{\lycee@bandeau{\signet{1}{#1}#1}}

\newcounter{exo}
\newlength{\espaceexo}\setlength{\espaceexo}{7pt}
\newlength{\espacetitre}\setlength{\espacetitre}{2pt}
\newcommand{\exo}[1][\relax]{%
  \par
  \ifapresbandeau\apresbandeaufalse\else\penalty-150\addvspace{\espaceexo}\fi
  \refstepcounter{exo}%
  \noindent\signet[exercice-\arabic{exo}]{2}{Exercice \arabic{exo}}%
  \textbf{\color{coulExo}\underline{Exercice \theexo}}%
  \ifx\relax#1\relax\else\textbf{ : #1}\fi
  \par\nopagebreak\vspace{\espacetitre}\nopagebreak
  \ignorespaces}

\setlist[enumerate]{nosep,leftmargin=*,itemsep=2pt,topsep=2pt}
\setlist[itemize]{nosep,leftmargin=*,topsep=2pt}
\newcommand{\choix}[3]{\par\hspace*{1em}%
  \makebox[0.3\linewidth][l]{a) #1}\makebox[0.3\linewidth][l]{b) #2}c) #3}
\newcommand{\deux}[2]{\par\hspace*{0.6em}%
  \makebox[0.48\linewidth][l]{#1}#2}
\newcommand{\trois}[3]{\par\hspace*{0.6em}%
  \makebox[0.32\linewidth][l]{#1}\makebox[0.32\linewidth][l]{#2}#3}

  \renewcommand{\exo}[1][\relax]{%
    \ifapresbandeau\apresbandeaufalse\else\par\penalty-150\fi
    \refstepcounter{exo}%
    \vspace{0.05cm}%
    \noindent\signet[exercice-\arabic{exo}]{2}{Exercice \arabic{exo}}%
    \textbf{\color{coulExo}\underline{Exercice \theexo}}%
    \ifx\relax#1\relax\else\textbf{ : #1}\fi%
    \par\nopagebreak
    \ignorespaces}
  \newcommand{\rappel}[1]{%
    \par\vspace{0.06cm}\nopagebreak
    \noindent\fcolorbox{black!35}{black!5}{%
      \parbox{\dimexpr\linewidth-2\fboxsep-2\fboxrule\relax}{%
        \small\textbf{Rappel.} #1}}%
    \par\vspace{0.06cm}\nopagebreak}
  \setlist[enumerate]{nosep,leftmargin=*}
  \setlist[itemize]{nosep,leftmargin=*}
  \setlength{\multicolsep}{3pt plus 1pt minus 1pt}
  \setlength{\premulticols}{10pt}
  \setlength{\postmulticols}{10pt}
  \newlist{expr}{enumerate}{1}
  \setlist[expr]{label={\Alph*\ $=$},nosep,leftmargin=*,itemsep=0pt}

\renewenvironment{center}{\par\addvspace{3pt}\centering}{\par\addvspace{3pt}}
\formulesserrees
\AtBeginDocument{\formulesserrees}
\clubpenalty=10000
\widowpenalty=10000
\displaywidowpenalty=10000
\brokenpenalty=10000
\setlength{\parskip}{0pt}

\RequirePackage{ragged2e}
\setlength{\RaggedRightRightskip}{0pt plus 1fil}
\AtBeginDocument{\RaggedRight\binoppenalty=10000 \relpenalty=10000 }
\g@addto@macro\@parboxrestore{\RaggedRight}
\makeatother

\newcommand{\venn}[6]{%
  \draw (-1.75,-1.35) rectangle (2.75,1.35);
  \node[font=\footnotesize] at (-1.5,1.1) {$\Omega$};
  \draw[coulCourbe,thick] (0,0) circle (1.1);
  \draw[coulCourbeB,thick] (1,0) circle (1.1);
  \node[coulCourbe] at (-0.95,0.95) {$#1$};
  \node[coulCourbeB] at (1.95,0.95) {$#2$};
  \node at (-0.6,0) {#3};
  \node at (0.5,0) {#4};
  \node at (1.6,0) {#5};
  \node at (2.3,-1.05) {#6};}
\tikzset{evt/.style={inner sep=2pt},
  proba/.style={font=\footnotesize,fill=white,inner sep=1pt}}
\newcommand{\arbre}[8]{%
  \node (O) at (0,0) {};
  \node[evt] (A) at (2.6,1.4) {$#1$};
  \node[evt] (nA) at (2.6,-1.4) {$\overline{#1}$};
  \node[evt] (AB) at (5.2,2.1) {$#2$};
  \node[evt] (AnB) at (5.2,0.7) {$\overline{#2}$};
  \node[evt] (nAB) at (5.2,-0.7) {$#2$};
  \node[evt] (nAnB) at (5.2,-2.1) {$\overline{#2}$};
  \draw (O.center) -- node[proba]{$#3$} (A);
  \draw (O.center) -- node[proba]{$#4$} (nA);
  \draw (A) -- node[proba]{$#5$} (AB);
  \draw (A) -- node[proba]{$#6$} (AnB);
  \draw (nA) -- node[proba]{$#7$} (nAB);
  \draw (nA) -- node[proba]{$#8$} (nAnB);}

\newcommand{\qcm}[4]{\par\hspace*{0.6em}\makebox[0.48\linewidth][l]{a) #1}b) #2
  \par\hspace*{0.6em}\makebox[0.48\linewidth][l]{c) #3}d) #4}
\newcommand{\qcmligne}[4]{\par\hspace*{0.6em}\makebox[0.24\linewidth][l]{a) #1}%
  \makebox[0.24\linewidth][l]{b) #2}\makebox[0.24\linewidth][l]{c) #3}d) #4}

\begin{document}

\titrefiche{1\textsuperscript{re} STMG --- Probabilités conditionnelles}{Fiche d'exercices du chapitre}

\begin{multicols}{2}\raggedcolumns

\partiefiche{Vocabulaire des événements}

\exo[Décrire des événements]
On tire au hasard un jeton dans un sac de $20$ jetons numérotés de $1$ à $20$. On note $A$
\og obtenir un multiple de $4$\fg{}, $B$ \og obtenir un nombre supérieur ou égal à
$15$\fg{} et $C$ \og obtenir un nombre impair\fg{}.
\begin{enumerate}[label=\arabic*)]
  \item Écrire l'univers $\Omega$, puis les événements $A$ et $B$ sous forme d'ensembles.
  \item Décrire par une phrase, puis écrire sous forme d'ensembles, les événements
        $A\cap B$ et $A\cup B$.
  \item Écrire l'événement $\overline{B}$ sous forme d'ensemble, puis le décrire par une
        phrase.
  \item Les événements $A$ et $C$ sont-ils incompatibles ? Et les événements $B$ et $C$ ?
\end{enumerate}

\exo[Écrire l'univers]
Pour chaque expérience aléatoire, écrire l'univers $\Omega$.
\begin{enumerate}[label=\alph*)]
  \item On lance une pièce et on note le côté visible : Pile (P) ou Face (F).
  \item On fait tourner une roue partagée en quatre secteurs : rouge (R), vert (V), bleu
        (B) et jaune (J).
  \item On forme un nombre de deux chiffres, chacun égal à $1$ ou à $2$.
  \item On lance un dé dont une face porte un $6$, deux faces un $1$ et trois faces un
        $2$, et on note le nombre obtenu.
\end{enumerate}

\exo[En phrase complète]
On choisit au hasard un élève d'un lycée. On note $A$ \og l'élève est en première\fg{} et
$B$ \og l'élève fait partie de l'association sportive\fg{}.
\begin{enumerate}[label=\arabic*)]
  \item Décrire par une phrase chacun des événements $A\cap B$, $A\cup B$, $\overline{A}$
        et $\overline{A}\cap B$.
  \item L'élève choisi est en seconde et fait partie de l'association sportive. Lesquels
        des événements $A$, $B$, $A\cap B$, $A\cup B$ et $\overline{A}$ sont réalisés ?
        Justifier par une phrase.
\end{enumerate}

\partiefiche{Diagramme de Venn}

\exo[Compléter un diagramme de Venn]
Une association de quartier compte $60$ bénévoles : $35$ réparent des vélos (ensemble
$R$), $28$ donnent de l'aide aux devoirs (ensemble $D$) et $12$ font les deux ; les autres
tiennent l'accueil.
\begin{enumerate}[label=\arabic*)]
  \item Représenter la situation par un diagramme de Venn.
  \item Combien de bénévoles font au moins une des deux activités $R$ et $D$ ? Aucune ?
        Une seule ?
\end{enumerate}

\exo[Partir de la zone commune]
Dans une classe de $30$ élèves, on note $S$ l'ensemble des élèves qui pratiquent un sport
en club et $M$ l'ensemble de ceux qui jouent d'un instrument de musique. $18$ élèves
pratiquent un sport en club, $8$ pratiquent un sport en club et jouent d'un instrument,
$5$ ne font ni l'un ni l'autre.
\begin{enumerate}[label=\arabic*)]
  \item Représenter la situation par un diagramme de Venn.
  \item Combien d'élèves jouent d'un instrument de musique ?
  \item Combien d'élèves font une seule des deux activités ?
\end{enumerate}

\exo[Lire un diagramme de Venn]
\begin{minipage}[c]{0.55\linewidth}\raggedright
Lors d'un salon de la mobilité, $200$ visiteurs ont pu essayer un vélo électrique
(ensemble $V$) et une trottinette électrique (ensemble $T$). Le diagramme ci-contre donne
les effectifs.
\end{minipage}\hfill
\begin{minipage}[c]{0.42\linewidth}\centering
\begin{tikzpicture}[scale=0.72]
  \venn{V}{T}{$70$}{$30$}{$40$}{$60$}
\end{tikzpicture}
\end{minipage}
\par\smallskip
\begin{enumerate}[label=\arabic*)]
  \item Combien de visiteurs ont essayé un vélo électrique ?
  \item Dire par une phrase ce que représente le nombre $30$, puis le nombre $60$.
  \item Combien de visiteurs ont essayé un seul des deux engins ?
  \item Décrire par une phrase l'ensemble $V\cup T$, puis donner son effectif.
\end{enumerate}

\partiefiche{Équiprobabilité}

\exo[Calculer des probabilités]
\begin{enumerate}[label=\alph*)]
  \item Une urne contient $30$ boules indiscernables, numérotées de $1$ à $30$. On en tire
        une au hasard. On note $A$ \og obtenir un multiple de $3$\fg{} et $B$ \og obtenir
        un multiple de $5$\fg{}. Calculer $P(A)$, $P(B)$ et $P(A\cap B)$.
  \item \begin{minipage}[t]{0.55\linewidth}\raggedright
        Le diagramme ci-contre décrit les $60$ bénévoles d'une association ($R$ :
        réparation de vélos ; $D$ : aide aux devoirs). On choisit un bénévole au hasard.
        Calculer $P(R\cap D)$, $P(R\cup D)$ et $P(\overline{R}\cap\overline{D})$.
        \end{minipage}\hfill
        \begin{minipage}[t]{0.4\linewidth}\centering\vspace{0pt}
        \begin{tikzpicture}[scale=0.7]
          \venn{R}{D}{$23$}{$12$}{$16$}{$9$}
        \end{tikzpicture}
        \end{minipage}
\end{enumerate}

\exo[Le stock d'une boutique]
Une boutique a en stock $100$ vestes, rangées selon leur couleur et leur taille.
\begin{center}
\renewcommand{\arraystretch}{1.3}
\begin{tabular}{|l|c|c|c|}
\hline
 & S & M & L\\
\hline
Bleue & $35$ & $15$ & $5$\\
\hline
Noire & $10$ & $20$ & $15$\\
\hline
\end{tabular}
\end{center}
On choisit une veste au hasard. Calculer la probabilité que la veste :
\begin{enumerate}[label=\arabic*)]
  \item soit noire ;
  \item soit de taille M ou L ;
  \item soit bleue et de taille S ;
  \item soit de taille L.
\end{enumerate}

\exo[Vrai ou faux]
Dire si chaque affirmation est vraie ou fausse, puis justifier par une phrase.
\begin{enumerate}[label=\alph*)]
  \item \og Avec un dé équilibré à six faces, la probabilité d'obtenir $7$ est égale à
        $0$.\fg{}
  \item \og Une probabilité peut être égale à $1{,}2$.\fg{}
  \item \og Une urne contient $5$ boules rouges et $3$ boules vertes, indiscernables ; la
        probabilité de tirer une boule verte est $\dfrac{3}{5}$.\fg{}
  \item \og On lance un dé équilibré dont une face porte un $6$, deux faces un $1$ et trois
        faces un $2$ : les trois issues possibles ont la même probabilité.\fg{}
\end{enumerate}

\partiefiche{Loi de probabilité}

\exo[Utiliser une loi de probabilité]
Un archer tire une flèche sur une cible. Le tableau donne la loi de probabilité du nombre
de points marqués.
\begin{center}
\renewcommand{\arraystretch}{1.3}\setlength{\tabcolsep}{3.5pt}
\begin{tabular}{|l|c|c|c|c|c|c|}
\hline
Points & $0$ & $2$ & $4$ & $6$ & $8$ & $10$\\
\hline
Probabilité & $0{,}05$ & $0{,}10$ & $0{,}15$ & $p$ & $0{,}30$ & $0{,}15$\\
\hline
\end{tabular}
\end{center}
\begin{enumerate}[label=\arabic*)]
  \item Calculer $p$.
  \item Calculer la probabilité de l'événement $A$ \og marquer au moins $6$ points\fg{}.
  \item Calculer la probabilité de l'événement $B$ \og marquer $4$ points ou moins\fg{}.
\end{enumerate}

\exo[À la caisse]
\begin{minipage}[c]{0.46\linewidth}\raggedright
Dans un supermarché, on a relevé pendant un an le temps d'attente en caisse. On choisit
un client au hasard ; le tableau donne la loi de probabilité de son temps d'attente.
\end{minipage}\hfill
\begin{minipage}[c]{0.52\linewidth}\centering
\renewcommand{\arraystretch}{1.3}\setlength{\tabcolsep}{3pt}
\begin{tabular}{|l|c|}
\hline
Attente & Probabilité\\
\hline
moins de $5$ min & $0{,}35$\\
\hline
de $5$ à $10$ min & $0{,}40$\\
\hline
de $10$ à $15$ min & $0{,}20$\\
\hline
plus de $15$ min & $p$\\
\hline
\end{tabular}
\end{minipage}
\begin{enumerate}[label=\arabic*)]
  \item Calculer $p$.
  \item Calculer la probabilité d'attendre moins de $10$ minutes.
  \item Calculer la probabilité d'attendre plus de $5$ minutes.
  \item Décrire par une phrase l'événement contraire de \og attendre moins de $10$
        minutes\fg{}, puis donner sa probabilité.
\end{enumerate}

\exo[Un dé truqué]
On lance un dé à six faces truqué : la face $1$ sort avec une probabilité de $0{,}05$, la
face $6$ avec une probabilité de $0{,}35$, et les quatre autres faces ont toutes la même
probabilité $p$.
\begin{enumerate}[label=\arabic*)]
  \item Calculer $p$.
  \item Calculer la probabilité de l'événement $A$ \og obtenir un nombre impair\fg{}.
  \item Calculer la probabilité de l'événement $B$ \og obtenir au moins $3$\fg{}.
  \item Quelle serait la probabilité de $A$ avec un dé équilibré ?
\end{enumerate}

\partiefiche{Événement contraire, réunion}

\exo[Contraire et réunion]
Un entrepôt expédie des colis. On choisit un colis au hasard et on note $A$ \og le colis
est volumineux\fg{}, $B$ \og le colis est fragile\fg{}, $C$ \og le colis est une
enveloppe\fg{}. On sait que $P(A)=0{,}25$ ; $P(B)=0{,}15$ ; $P(A\cap B)=0{,}05$ ;
$P(C)=0{,}3$ ; et une enveloppe n'est jamais volumineuse.
\begin{enumerate}[label=\arabic*)]
  \item Calculer $P(\overline{A})$ et traduire le résultat par une phrase.
  \item Calculer la probabilité qu'un colis soit volumineux ou fragile.
  \item En déduire la probabilité qu'un colis ne soit ni volumineux ni fragile.
  \item Calculer $P(A\cup C)$.
\end{enumerate}

\exo[Retrouver une intersection]
Pour trois événements $A$, $B$ et $C$ d'une même expérience, on sait que $P(A)=0{,}2$ ;
$P(B)=0{,}35$ ; $P(A\cup B)=0{,}45$ ; $P(C)=0{,}4$ ; et $B$ et $C$ sont incompatibles.
\begin{enumerate}[label=\arabic*)]
  \item Calculer $P(\overline{B})$.
  \item Calculer $P(B\cup C)$.
  \item Calculer $P(A\cap B)$.
  \item Calculer $P(\overline{A}\cap\overline{B})$.
\end{enumerate}

\exo[Au cinéma]
Ce mois-ci, un cinéma a vendu $5\,000$ billets : $30\,\%$ avec un menu pop-corn
(ensemble $A$), $50\,\%$ au tarif réduit (ensemble $B$) et $10\,\%$ avec un menu pop-corn
au tarif réduit.
\begin{enumerate}[label=\arabic*)]
  \item Calculer le nombre de billets de $A$, de $B$ et de $A\cap B$, puis représenter les
        ventes par un diagramme de Venn.
  \item On choisit un billet vendu au hasard. Calculer la probabilité qu'il ait un menu
        pop-corn ou qu'il soit au tarif réduit.
  \item Calculer la probabilité qu'il n'ait ni menu pop-corn ni tarif réduit.
  \item Calculer la probabilité qu'il soit dans un seul des deux cas.
\end{enumerate}

\partiefiche{Arbre des possibles}

\exo[Dénombrer avec un arbre]
Le local à vélos d'un lycée s'ouvre avec un code de deux chiffres, chacun égal à $1$, à
$2$, à $3$ ou à $4$. On essaie un code au hasard.
\begin{enumerate}[label=\arabic*)]
  \item Représenter toutes les issues par un arbre. Combien y en a-t-il ?
  \item Calculer la probabilité de trouver le bon code.
  \item Calculer la probabilité de l'événement $A$ \og les deux chiffres sont égaux\fg{}.
\end{enumerate}
Un sac contient trois jetons : un rouge (R), un vert (V) et un jaune (J). On tire un jeton
au hasard, on le garde, puis on en tire un deuxième.
\begin{enumerate}[label=\arabic*),start=4]
  \item Représenter toutes les issues par un arbre. Combien y en a-t-il ?
  \item Calculer la probabilité de l'événement $B$ \og le premier jeton est rouge\fg{}.
  \item Calculer la probabilité de l'événement $C$ \og le jeton rouge fait partie des deux
        jetons tirés\fg{}.
\end{enumerate}

\exo[Trois lancers]
On lance trois fois de suite une pièce équilibrée ; on note P pour Pile et F pour Face.
\begin{enumerate}[label=\arabic*)]
  \item Représenter toutes les issues par un arbre. Combien y en a-t-il ?
  \item Calculer la probabilité d'obtenir trois fois Pile.
  \item Calculer la probabilité de l'événement $A$ \og obtenir exactement une fois
        Pile\fg{}.
  \item À l'aide de l'événement contraire, calculer la probabilité d'obtenir au moins une
        fois Pile.
\end{enumerate}

\exo[La tenue du cycliste]
Pour une sortie, un cycliste choisit au hasard un maillot parmi trois (rouge, bleu, blanc)
et un casque parmi deux (noir, blanc).
\begin{enumerate}[label=\arabic*)]
  \item Représenter toutes les tenues possibles par un arbre. Combien y en a-t-il ?
  \item Calculer la probabilité que le maillot et le casque soient blancs.
  \item On note $M$ \og le maillot est blanc\fg{} et $C$ \og le casque est blanc\fg{}.
        Calculer $P(M)$ et $P(C)$, puis $P(M\cup C)$ avec la formule de la réunion.
\end{enumerate}

\partiefiche{Probabilité conditionnelle et tableau croisé}

\exo[Probabilités dans un tableau croisé]
Une piscine compte $400$ abonnés. Le tableau, à compléter, les range selon leur formule et
le moment où ils viennent le plus souvent.
\begin{center}
\renewcommand{\arraystretch}{1.3}
\begin{tabular}{|l|c|c|c|}
\hline
 & Mensuelle & Annuelle & Total\\
\hline
Matin & $100$ & $100$ & \\
\hline
Soir & $50$ & & \\
\hline
Total & & & $400$\\
\hline
\end{tabular}
\end{center}
On choisit un abonné au hasard ; on note $A$ \og l'abonné a une formule annuelle\fg{} et
$S$ \og l'abonné vient le soir\fg{}.
\begin{enumerate}[label=\arabic*)]
  \item Compléter le tableau.
  \item Calculer $P(A)$ et $P(A\cap S)$.
  \item Calculer $P_A(S)$ et traduire le résultat par une phrase.
  \item Calculer $P_S(A)$ et traduire le résultat par une phrase.
\end{enumerate}

\exo[Une équipe de rugby]
Une équipe de rugby compte $40$ joueurs : $24$ avants et $16$ arrières. $18$ avants
mesurent plus de $1{,}85$ m, contre $2$ arrières. On note $A$ \og le joueur est un
avant\fg{} et $G$ \og le joueur mesure plus de $1{,}85$ m\fg{}.
\begin{enumerate}[label=\arabic*)]
  \item Présenter ces données dans un tableau croisé d'effectifs.
  \item On choisit un joueur au hasard. Calculer $P(A)$, $P(G)$ et $P(A\cap G)$.
  \item Calculer $P_A(G)$, puis traduire ce résultat par une phrase.
  \item Calculer $P_G(A)$, puis traduire ce résultat par une phrase.
  \item Expliquer par une phrase la différence entre $P(A\cap G)$ et $P_A(G)$.
\end{enumerate}

\exo[Au magasin de cycles]
Un magasin de cycles a vendu $400$ vélos : $55\,\%$ de VTT, $30\,\%$ de VTC, et les autres
sont des vélos de ville. $60\,\%$ des VTT, $50\,\%$ des VTC et $80\,\%$ des vélos de ville
vendus sont à assistance électrique.
\begin{enumerate}[label=\arabic*)]
  \item Compléter le tableau d'effectifs.
        \begin{center}
        \renewcommand{\arraystretch}{1.3}
        \begin{tabular}{|l|c|c|c|}
        \hline
         & Électrique & Non électrique & Total\\
        \hline
        VTT & & & \\
        \hline
        VTC & & & \\
        \hline
        Ville & & & \\
        \hline
        Total & & & $400$\\
        \hline
        \end{tabular}
        \end{center}
  \item On choisit un vélo vendu au hasard ; on note $E$ \og le vélo est électrique\fg{} et
        $T$ \og le vélo est un VTT\fg{}. Calculer $P(E)$, puis $P(\overline{E})$.
  \item Calculer $P_E(T)$ et traduire le résultat par une phrase.
  \item Un client a acheté un vélo non électrique. Calculer la probabilité que ce soit un
        VTC.
\end{enumerate}

\partiefiche{Formule des probabilités conditionnelles}

\exo[Traduire et utiliser la formule]
\begin{enumerate}[label=\alph*)]
  \item Dans une ville, $40\,\%$ des maisons ont une pompe à chaleur, et $10\,\%$ des maisons
        ont une pompe à chaleur et des panneaux solaires. On choisit une maison au hasard ;
        on note $C$ \og la maison a une pompe à chaleur\fg{} et $S$ \og la maison a des
        panneaux solaires\fg{}. Traduire les données par des probabilités, puis calculer
        $P_C(S)$ et traduire le résultat par une phrase.
  \item Dans un lycée, $60\,\%$ des élèves déjeunent à la cantine ; parmi eux, $30\,\%$
        choisissent le menu végétarien. On choisit un élève au hasard ; on note $D$
        \og l'élève déjeune à la cantine\fg{} et $V$ \og l'élève choisit le menu
        végétarien\fg{}. Traduire les données par des probabilités, puis calculer
        $P(D\cap V)$ et traduire le résultat par une phrase.
\end{enumerate}

\exo[Conditionnelle ou pas ?]
Une station météo choisit un jour de l'année au hasard. On note $R$ \og il a plu ce
jour-là\fg{} et $V$ \og le vent a soufflé fort ce jour-là\fg{}.
\begin{enumerate}[label=\arabic*)]
  \item Pour chaque information, dire s'il s'agit d'une probabilité conditionnelle, puis
        l'écrire avec la notation qui convient.
        \begin{enumerate}[label=\alph*)]
          \item Dans l'année, $40\,\%$ des jours sont pluvieux.
          \item $75\,\%$ des jours pluvieux sont ventés.
          \item Parmi les jours sans vent fort, $20\,\%$ sont pluvieux.
          \item $30\,\%$ des jours sont à la fois pluvieux et ventés.
          \item Quand le vent souffle fort, il pleut $6$ jours sur $10$.
        \end{enumerate}
  \item Vérifier que les informations a), b) et d) sont cohérentes.
  \item Calculer la probabilité qu'un jour choisi au hasard ne soit pas pluvieux.
\end{enumerate}

\exo[Au club nautique]
Dans un club nautique, $60\,\%$ des adhérents font de la voile, $50\,\%$ du kayak, et
$15\,\%$ font les deux. On choisit un adhérent au hasard ; on note $V$ \og l'adhérent fait
de la voile\fg{} et $K$ \og l'adhérent fait du kayak\fg{}.
\begin{enumerate}[label=\arabic*)]
  \item Donner $P(V)$, $P(K)$ et $P(V\cap K)$.
  \item Calculer $P_V(K)$ et traduire le résultat par une phrase.
  \item Calculer la probabilité qu'un adhérent qui fait du kayak fasse aussi de la voile.
  \item Calculer la probabilité qu'un adhérent fasse au moins une des deux activités.
\end{enumerate}

\partiefiche{Arbre pondéré}

\exo[Construire un arbre pondéré]
Dans une salle d'escalade, $40\,\%$ des grimpeurs sont débutants. Parmi les débutants,
$30\,\%$ réussissent la voie rouge ; parmi les autres grimpeurs, $80\,\%$. On choisit un
grimpeur au hasard ; on note $D$ \og le grimpeur est débutant\fg{} et $R$ \og le grimpeur
réussit la voie rouge\fg{}.
\begin{enumerate}[label=\arabic*)]
  \item Construire l'arbre pondéré de la situation.
  \item Calculer la probabilité qu'un grimpeur soit débutant et réussisse la voie rouge ;
        traduire le résultat par une phrase.
  \item Calculer la probabilité qu'un grimpeur soit débutant et échoue sur la voie rouge.
\end{enumerate}

\exo[Le jeu de la fête de quartier]
Lors d'une fête de quartier, un jeu se déroule en deux étapes. On tire au hasard une carte
parmi $40$ cartes numérotées de $1$ à $40$. Si le numéro est inférieur ou égal à $10$, on
fait tourner une roue de $10$ secteurs égaux dont $6$ portent une étoile ; sinon, une autre
roue de $10$ secteurs égaux dont $3$ portent une étoile. On gagne un lot si la roue
s'arrête sur une étoile. On note $N$ \og le numéro est inférieur ou égal à $10$\fg{} et $E$
\og la roue s'arrête sur une étoile\fg{}.
\begin{enumerate}[label=\arabic*)]
  \item Justifier que $P(N)=0{,}25$ et que $P_N(E)=0{,}6$.
  \item Construire l'arbre pondéré de la situation.
  \item Calculer la probabilité de tirer un numéro inférieur ou égal à $10$ et de gagner un
        lot.
  \item Calculer $P(\overline{N}\cap E)$.
\end{enumerate}

\exo[Compléter et lire un arbre]
\begin{minipage}[c]{0.4\linewidth}\raggedright
À la sortie d'un lycée, on choisit un élève au hasard. On note $B$ \og l'élève rentre en
bus\fg{} et $H$ \og l'élève est chez lui avant 17~h~30\fg{}.
\end{minipage}\hfill
\begin{minipage}[c]{0.57\linewidth}\centering
\begin{tikzpicture}[scale=0.78]
  \arbre{B}{H}{0{,}65}{\dots}{0{,}4}{\dots}{0{,}7}{\dots}
\end{tikzpicture}
\end{minipage}
\begin{enumerate}[label=\arabic*)]
  \item Compléter l'arbre.
  \item Traduire par une phrase les probabilités $P(B)=0{,}65$ et
        $P_{\overline{B}}(H)=0{,}7$.
  \item Calculer la probabilité qu'un élève rentre en bus et soit chez lui après 17~h~30.
\end{enumerate}

\partiefiche{Formule des probabilités totales}

\exo[Probabilités totales]
\begin{minipage}[c]{0.4\linewidth}\raggedright
Dans une salle d'escalade, l'arbre ci-contre décrit les grimpeurs ($D$ \og le grimpeur est
débutant\fg{}, $R$ \og il réussit la voie rouge\fg{}).
\end{minipage}\hfill
\begin{minipage}[c]{0.57\linewidth}\centering
\begin{tikzpicture}[scale=0.78]
  \arbre{D}{R}{0{,}4}{0{,}6}{0{,}3}{0{,}7}{0{,}8}{0{,}2}
\end{tikzpicture}
\end{minipage}
\begin{enumerate}[label=\arabic*)]
  \item Calculer $P(R)$.
  \item Un grimpeur a réussi la voie rouge. Calculer la probabilité qu'il soit débutant.
\end{enumerate}
Un magasin d'électroménager achète $60\,\%$ de ses aspirateurs à un premier fournisseur,
et les autres à un second. Pendant la garantie, $2\,\%$ des aspirateurs du premier
fournisseur tombent en panne, contre $7\,\%$ de ceux du second. On choisit un aspirateur
vendu au hasard ; on note $F$ \og il vient du premier fournisseur\fg{} et $D$ \og il tombe
en panne pendant la garantie\fg{}.
\begin{enumerate}[label=\arabic*),start=3]
  \item Construire l'arbre pondéré, puis calculer $P(D)$.
  \item Un aspirateur est tombé en panne. Calculer la probabilité qu'il vienne du premier
        fournisseur.
\end{enumerate}

\exo[Les visiteurs d'un site]
Une boutique en ligne étudie les visiteurs de son site : $75\,\%$ arrivent par un moteur
de recherche, les autres par la lettre d'information de la boutique. $20\,\%$ des
visiteurs venus par un moteur de recherche passent commande, contre $40\,\%$ de ceux
venus par la lettre d'information. On choisit un visiteur au hasard ; on note $M$ \og le
visiteur arrive par un moteur de recherche\fg{} et $C$ \og le visiteur passe
commande\fg{}.
\begin{enumerate}[label=\arabic*)]
  \item Construire l'arbre pondéré de la situation.
  \item Calculer $P(M\cap C)$ et $P(\overline{M}\cap C)$.
  \item En déduire $P(C)$.
  \item Un visiteur a passé commande. Calculer la probabilité qu'il soit arrivé par un
        moteur de recherche.
\end{enumerate}

\exo[Aller en train, revenir en bus ?]
Pour une excursion, chaque voyageur fait l'aller en train ou en bus, puis le retour en
train ou en bus. À l'aller, $60\,\%$ des voyageurs prennent le train. Parmi ceux qui ont
pris le train à l'aller, $80\,\%$ le reprennent au retour ; parmi ceux qui ont pris le bus à
l'aller, $40\,\%$ prennent le train au retour. On note $A$ \og l'aller se fait en
train\fg{} et $R$ \og le retour se fait en train\fg{}.
\begin{enumerate}[label=\arabic*)]
  \item Construire l'arbre pondéré de la situation.
  \item Calculer la probabilité qu'un voyageur fasse l'aller et le retour en train.
  \item Calculer $P(R)$.
  \item Calculer la probabilité qu'un voyageur utilise les deux moyens de transport.
  \item Un voyageur est revenu en train. Calculer la probabilité qu'il soit aussi allé en
        train.
\end{enumerate}

\partiefiche{Événements indépendants}

\exo[Indépendants ou pas ?]
\begin{enumerate}[label=\alph*)]
  \item Un atelier a révisé $250$ vélos électriques. On choisit un vélo au hasard ; on note
        $F$ \og les freins sont à régler\fg{} et $N$ \og un pneu est à changer\fg{}. Les
        événements $F$ et $N$ sont-ils indépendants ?
        \begin{center}
        \renewcommand{\arraystretch}{1.3}
        \begin{tabular}{|l|c|c|c|}
        \hline
         & $F$ & $\overline{F}$ & Total\\
        \hline
        $N$ & $20$ & $30$ & $50$\\
        \hline
        $\overline{N}$ & $80$ & $120$ & $200$\\
        \hline
        Total & $100$ & $150$ & $250$\\
        \hline
        \end{tabular}
        \end{center}
  \item Dans un lycée, un jour de pluie, $30\,\%$ des élèves viennent en bus ; sur
        l'ensemble de l'année, $20\,\%$ des élèves viennent en bus. On choisit un jour et un
        élève au hasard ; on note $R$ \og il pleut ce jour-là\fg{} et $B$ \og l'élève vient
        en bus\fg{}. Les événements $R$ et $B$ sont-ils indépendants ?
  \item Une agence bancaire a deux distributeurs de billets, qui tombent en panne
        indépendamment l'un de l'autre. Un jour donné, le premier fonctionne toute la
        journée avec une probabilité de $0{,}98$ et le second avec une probabilité de
        $0{,}9$. Calculer la probabilité que les deux distributeurs fonctionnent toute la
        journée, puis la probabilité qu'au moins un tombe en panne.
\end{enumerate}

\exo[Une pièce et un dé]
On lance une pièce équilibrée, puis un dé équilibré à six faces. On note $A$ \og la pièce
donne Pile\fg{} et $B$ \og le dé donne $6$\fg{}.
\begin{enumerate}[label=\arabic*)]
  \item Représenter toutes les issues par un arbre. Combien y en a-t-il ?
  \item Calculer $P(A)$, $P(B)$ et $P(A\cap B)$.
  \item Calculer $P_A(B)$. Les événements $A$ et $B$ sont-ils indépendants ?
  \item Vérifier que $P(A\cap B)=P(A)\times P(B)$.
\end{enumerate}

\exo[Une commande en retard]
Pendant les soldes, une commande d'une boutique en ligne part en retard si un article
est en rupture de stock (événement $A$) ou si le transporteur ne peut pas la prendre le
jour même (événement $B$). Les événements $A$ et $B$ sont indépendants, avec
$P(A)=0{,}05$ et $P(B)=0{,}02$.
\begin{enumerate}[label=\arabic*)]
  \item Calculer la probabilité qu'un article soit en rupture de stock et que le
        transporteur ne puisse pas prendre la commande.
  \item Calculer la probabilité que la commande parte en retard.
  \item En déduire la probabilité qu'elle parte à temps.
\end{enumerate}

\partiefichesansnum{Exercices type épreuve anticipée}

\exo[QCM]
Pour chaque question, une seule réponse est exacte.
\begin{enumerate}[label=\arabic*.]
  \item On lance un dé truqué à six faces : la face $1$ sort avec une probabilité de
        $0{,}1$ et la face $2$ avec une probabilité de $0{,}25$. La probabilité d'obtenir
        au moins $3$ est :
        \qcmligne{$0{,}35$}{$0{,}75$}{$0{,}5$}{$0{,}65$}
  \item Une urne contient $25$ jetons indiscernables, dont $10$ gagnants. On tire un jeton
        au hasard. La probabilité qu'il soit gagnant est :
        \qcmligne{$0{,}1$}{$0{,}25$}{$0{,}4$}{$2{,}5$}
  \item Un magasin a vendu $200$ vélos : $130$ pour adulte, dont $90$ électriques, et
        $70$ pour enfant, dont $10$ électriques. On choisit au hasard un client parmi ceux
        qui ont acheté un vélo pour enfant. La probabilité que ce vélo soit électrique
        est :
        \qcmligne{$\dfrac{10}{70}$}{$\dfrac{10}{200}$}{$\dfrac{10}{100}$}{$\dfrac{70}{200}$}
  \item $A$ et $B$ sont deux événements tels que $P(A)=0{,}4$ et $P_A(B)=0{,}5$. Alors
        $P(A\cap B)$ vaut :
        \qcmligne{$0{,}9$}{$0{,}1$}{$0{,}2$}{$0{,}8$}
  \item Dans une entreprise, parmi les salariés qui viennent en train, $30\,\%$ ont un
        abonnement annuel. On choisit un salarié au hasard ; on note $T$ \og il vient en
        train\fg{} et $A$ \og il a un abonnement annuel\fg{}. Cette information s'écrit :
        \qcm{$P(T\cap A)=0{,}3$}{$P_T(A)=0{,}3$}{$P_A(T)=0{,}3$}{$P(A)=0{,}3$}
  \item $A$ et $B$ sont deux événements indépendants, avec $P(A)=0{,}5$ et $P(B)=0{,}4$.
        Alors $P_A(B)$ vaut :
        \qcmligne{$0{,}5$}{$0{,}2$}{$0{,}9$}{$0{,}4$}
\end{enumerate}

\exo[Vrai ou faux ?]
Une salle de sport compte $400$ abonnés, rangés selon leur formule et leur
participation aux cours collectifs.
\begin{center}
\renewcommand{\arraystretch}{1.3}
\begin{tabular}{|l|c|c|c|}
\hline
 & Avec cours & Sans cours & Total\\
\hline
Annuelle & $60$ & $180$ & $240$\\
\hline
Mensuelle & $40$ & $120$ & $160$\\
\hline
Total & $100$ & $300$ & $400$\\
\hline
\end{tabular}
\end{center}
On choisit un abonné au hasard ; on note $A$ \og il a une formule annuelle\fg{} et $C$
\og il suit des cours collectifs\fg{}. Pour chacune des quatre affirmations, dire si elle
est vraie ou fausse, en justifiant.
\begin{enumerate}[label=\arabic*.]
  \item \emph{Au moins $50\,\%$ des abonnés ont une formule annuelle.}
  \item \emph{La probabilité qu'un abonné n'ait pas de formule annuelle et ne suive pas
        de cours collectifs est égale à $0{,}25$.}
  \item \emph{Au moins un tiers des abonnés à formule mensuelle suivent des cours
        collectifs.}
  \item \emph{Les événements $A$ et $C$ sont indépendants.}
\end{enumerate}

\exo[Le filtre anti-spam]
Une entreprise teste le filtre anti-spam de sa messagerie sur $500$ courriels reçus.
\begin{center}
\renewcommand{\arraystretch}{1.3}
\begin{tabular}{|l|c|c|c|}
\hline
 & Bloqués & Non bloqués & Total\\
\hline
Spams & $96$ & $4$ & $100$\\
\hline
Autres & $24$ & $376$ & $400$\\
\hline
Total & $120$ & $380$ & $500$\\
\hline
\end{tabular}
\end{center}
On choisit un de ces courriels au hasard ; on note $S$ \og le courriel est un spam\fg{} et
$B$ \og le filtre bloque le courriel\fg{}. Les résultats sont donnés sous forme décimale.
\begin{enumerate}[label=\arabic*.]
  \item Montrer que $P(S)=0{,}2$.
  \item Décrire par une phrase l'événement $S\cap B$, puis calculer sa probabilité.
  \item Calculer $P_S(B)$ et traduire le résultat par une phrase.
  \item Le filtre est jugé fiable si la probabilité qu'un courriel ne soit pas un spam,
        sachant qu'il est bloqué, est inférieure ou égale à $0{,}15$. Le filtre est-il
        fiable ? On donne $120=5\times 24$.
  \item Calculer la probabilité que le filtre se trompe (un spam non bloqué, ou un autre
        courriel bloqué).
\end{enumerate}

\end{multicols}

\end{document}
